\documentclass{article}
\usepackage{tikz}
\usetikzlibrary{angles, quotes}
\usepackage{amssymb}
\usepackage{pgfplots}
\pgfplotsset{width=10cm,compat=1.9}
\usepgfplotslibrary{statistics}
\usepackage[left=2.54cm,right=2.54cm,top=2.54cm,bottom=2.54cm]{geometry}
\usepackage{tasks} %Paquete necesario para la producción de items horizontales para algunos ejercicios de tarea empleando el comando /begin{multicols}}{#}.../end{multicols} para ayudarnos a reducir las páginas
\usepackage{pdflscape} %comando que permite cambiar a orientación horizontal las hojas%
\usepackage{soul} %Paquete para permitir la compilación de acentos usando el comando {\hl{}}}\\
\sethlcolor{yellow(munsell)} %comando para definir el color para el subrayado usando el comando \hl
\usepackage[utf8]{inputenc}
\usepackage[spanish]{babel}
\usepackage{icomma}
\usepackage{ragged2e}
\usepackage{siunitx}
\usepackage{url}
\usepackage[colorlinks=true, urlcolor=blue,  linkcolor=black, citecolor=green]{hyperref}
\usepackage{pdfpages}
\usepackage{blindtext} %Paquete que genera texto ficticio (dummy text) con el comando: \blindtext
\usepackage{cancel} %in the preamble gives you four different modes of striking through: \cancel{text to cancel} draws a diagonal line (slash) through its argument, \bcancel{text to cancel} uses the negative slope (a backslash), \xcancel{text to cancel} draws an X (actually \cancel plus \bcancel), \cancelto{〈value〉}{〈expression〉} draws a diagonal arrow through the 〈expression〉pointing to the 〈value〉 (math-mode only)
\usepackage{csquotes}
\usepackage{afterpage}
\usepackage{parskip} 
\usepackage{float}
\usepackage{enumitem}
\usepackage{multicol}%Paquete que permite la creación aislada de columnas de texto. De esta forma se reduce la cantidad de páginas en nuestro documento
\usepackage{lipsum} %Paquete que genera texto de relleno al igual que \usepackage{blindtext}, se usa con el comando \lipsum[#-#] los #(hashtags) delimitan el número de páginas que deseamos ocupar del paquete lipsum para usarlos como relleno. 

% Definir el comando personalizado para la nota
\newcommand{\mynote}[1]{%
\begin{tikzpicture}[baseline=-0.75ex]
\node[draw, circle, fill=Apple Green!60, inner sep=2pt] (note) {\textbf{Nota:}};
\end{tikzpicture}%
\ \textit{#1}%
}

\newenvironment{Figura}
{\par\medskip\noindent\minipage{\linewidth}}
{\endminipage\par\medskip}
\usepackage{caption}
\usepackage[
backend=biber,
sorting=none,
url=true
]{biblatex} %bibliografía
\addbibresource{biblio.bib}
% Establecer el espacio entre las entradas de la bibliografía
\setlength{\bibitemsep}{1\baselineskip} % Puedes ajustar el valor según tus preferencias
\usepackage{amsmath,amsthm,amssymb,amsfonts}
\usepackage{pifont} %Permite la compilación del comando \xmark: es el símbolo de la tachita
\usepackage{mathtools} %permite la compilación de símbolos para matrices
\usepackage{empheq} %Paquete que se relaciona con \usepackage[most]{tcolorbox} para la creación de cuadros/cajas de colores para delimitar los resultados matemáticos o líneas de texto usando la línea de comando: \begin{empheq}[box={\mymath[colback=orange(webcolor)!70,drop lifted shadow]}]{equation*} \end{empheq}
\usepackage[most]{tcolorbox} %Paquete que se relaciona con \usepackage{empheq} para la creación de cuadros/cajas de colores para delimitar los resultados matemáticos o líneas de texto
\renewcommand{\qedsymbol}{$\blacksquare$}
\usetikzlibrary{positioning,decorations.pathreplacing} %Paquete que permite la compilación de llaves para cuadros sinópticos (brace diagram) 
\usepackage{schemata} %The schemata package is designed just for these brace diagrams
\usepackage{booktabs}

\newcommand\AB[2]{\schema{\schemabox{#1}}{\schemabox{#2}}} %comando o paquete necesario para crear cuadros sinópticos

\newtcbox{\mymath}[1][]{%
nobeforeafter, math upper, tcbox raise base,
enhanced, colframe=blue!30!black,
colback=blue!30, boxrule=1pt,
#1} %Comando importante para encasillar los resultados matemáticos en cajas de diferentes colores, se relaciona con los paquetes \usepackage{empheq} y \usepackage[most]{tcolorbox} 

\newenvironment{sysmatrix}[1]
{\left(\begin{array}{@{}#1@{}}}
{\end{array}\right)}
\newcommand{\ro}[1]{%
\xrightarrow{\mathmakebox[\rowidth]{#1}}%
}
\newlength{\rowidth}% row operation width
\AtBeginDocument{\setlength{\rowidth}{3em}}  %Comando importante para laproducción de lineas de operaciones entre matrices, método de Gauss o Gauss Jordan se relaciona con \newenvironment{sysmatrix}

%formato para cambiar el horario a español
\usepackage[spanish]{datetime2}
\DTMsetdatestyle{spanish}

\renewcommand{\today}{\DTMdisplaydate{\the\year}{\the\month}{\the\day }{-1}}
%%%%%%%%%%%%%%%%%%%%%%%%%%%%%%%%%%%%%%%%%

\usepackage{datetime}
\newcommand{\mycurrenttime}{\xxivtime}

%%%%%%%%%%%%%%%%%%%%%%%%%%%%%%%%%%%%%%%%%%%%%%%%%%%%%%%%%%%%%%%%%%%%%%%%%%%%%%%%%%%%%%%%%%%%%%%%%%%%%%%%%%%%%%%%%%%%%%%%%%%%%%%%%%%%%%%%%%%%%%%%%


%%%%%%%%%%%%%Caja de color para el título%%%%%%%%%%%%%%%%%%%%%%%%%%%%%%%%%%%%%%%%%%%%%%%%%%%%%%

\definecolor{myframecolor}{RGB}{1, 45, 75} % Prussian Blue
\definecolor{myboxcolor}{RGB}{0, 107, 92} % Apple Green

\newtcolorbox{mybox}{
enhanced,
colback=myboxcolor!7, % Color de fondo del cuadro
colframe=myframecolor, % Color del marco del cuadro
arc=0pt,
boxrule=1pt,
borderline west={2mm}{-10mm}{myframecolor}, % Borde en el lado izquierdo
sharp corners=southwest,
width=\linewidth,
before=\par\vspace{\bigskipamount}, % Espacio antes del cuadro
after=\par\vspace{\bigskipamount} % Espacio después del cuadro
}


\newcommand{\euler}{e} %Comando para producir letra e de euler
\newenvironment{remark}{\par\vfill\footnotesize % Vertical white space above the remark and smaller font size
\begin{list}{}{
	\leftmargin=80pt % Indentation on the left
	\rightmargin=60pt}\item\ignorespaces % Indentation on the right
\makebox[-2.5pt]{\begin{tikzpicture}[overlay]
		\node[draw=Horizon!60,line width=2.5pt,circle,fill=Horizon!25,font=\sffamily\bfseries,inner sep=4pt,outer sep=0pt] at (-19pt,5pt){\textcolor{Horizon}{Nota}};\end{tikzpicture}} % Blue Nota in a circle
\advance\baselineskip -1pt}{\end{list}\vskip5pt} % Tighter line spacing and white space after remark
\usepackage{graphicx}

\usepackage{titling}

\input{colores.tex}

%%Producir signos de cita textual``Asignatura''

\newcommand{\xmark}{\ding{55}} %%%comando que genera la tachita

\newenvironment{MyColorPar}[1]{%
\leavevmode\color{#1}\ignorespaces%
}{%
}%

%%%%%%%%%%%%%%%%%%% Título de la tarea, Nombre de alumno

\title{Clase \# 40 - Mecánica Teórica}
\author{\emph{JA} $\boldsymbol{\mid}$}
\date{\today}

%%%%%%%%%%%%%%%%%%% Datos de la Materia

\usepackage{fancyhdr}
\fancypagestyle{plain}{%  the preset of fancyhdr 
\fancyhf{} % clear all header and footer fields
\fancyfoot[R]{
\includegraphics[width=2cm]{LogoFCFMBUAP (1).png}}
\fancyfoot[L]{{\bfseries{\thedate{}}} a las {\bfseries{\mycurrenttime{}}} horas{} (GMT-6, H. Puebla de Zaragoza, Pue)}
\fancyhead[L]{Mecánica teórica}
\fancyhead[R]{\theauthor}
}
\makeatletter
\def\@maketitle{%
\newpage
\null
\vskip 1em%
\begin{center}%
\let \footnote \thanks
{\LARGE \@title \par}%
\vskip 1em%
%{\large \@date}%
\end{center}%
\par
\vskip 1em}
\makeatother

\usepackage{lipsum}  
\usepackage{cmbright}


%%%%%%%%%%%%%%%%%%%%%%%%%%%%%%%%%%%%%%%%%%%%%%%%%%%%%%%%%%%%%%%%%%%%%%%%%%%%%%%%%%%%%%%%%%%%%%%%%%%%%%%%%%%%%%%%%%%%%%%%%%%%%%%%%%%%%%%%%%%%%%%%%%%%%%%%%%%%%%%%%%%%%%%%%%%%%%%%%%%%%%%%%%%%%%%%%%%%%%

\begin{document}



\maketitle


%%%%%%%%%%%%%%%%%%% Datos del profesor, campus y aula

\begin{mybox}
\noindent\begin{tabular}{@{}ll}
	{\bfseries{Profesor}} & Dr. Gilberto Silva Ortigoza\\
	\textcolor{prussianblue}{\bfseries{Campus}} / \textcolor{trueblue}{\bfseries{Facultad}}  & \textcolor{prussianblue}{\bfseries{C.U. BUAP}} /  \textcolor{trueblue}{\bfseries{FCFM}} \\
	\textcolor{prussianblue}{\bfseries{Edificio}} /  \textcolor{trueblue}{\bfseries{EMA4}}    & 	\textcolor{prussianblue}{\bfseries{FM9}} /  \textcolor{trueblue}{\bfseries{302}}   
\end{tabular} 
\end{mybox} \vspace{1cm}

\tableofcontents

\vspace{2cm}

\section{Clase del día de hoy}


\[ \delta \int\limits_{t_{1}}   \left[  \sum\limits_{j=1}^{n} p_{j}\dot{q}_{j} - H(q_{i},p_{i},t)   \right]dt=0 \]



Tenemos:


\[ \dot{q}_{i} = \frac{ \partial H }{\partial p_{i}} \]


\[ \dot{p}_{i} = - \frac{\partial H}{\partial_{i}}    \]



Tenemos:


\[ Q_{i} = Q_{i}(q_{j},p_{j},t) \]


\[P_{i} = P_{i}(q_{j},p_{j},t)\]


\[ J \neq 0 \]



Tenemos:



\[ \left[ Q_{i},Q_{j}  \right] =0  \]


\[ \left[ P_{i},P_{j}  \right] =0  \]


\[ \left[ Q_{i}, P_{j} \right] = \delta_{ij}   \]


T. canónica.


Tenemos:



\[ \dot{Q}_{i} = \frac{\partial K}{\partial P_{i}}   \]



\[ \partial{P}_{i}  = - \frac{ \partial K  }{\partial Q_{i}}  \]



Tenemos:




\[  \delta \int\limits \left[  \sum\limits_{j=1}^{n} P_{i}\dot{Q}_{i} - K(Q_{i},P_{i},t)   \right]  dt =0 \]



Tenemos:



\[ L^{\prime} = L + \frac{dF}{dt}(q,t) \]




Tenemos:





\[ \left[  \sum\limits_{j=1}^{n} p_{j}\dot{q}_{j} - H   \right] dt + \frac{dF}{dt} dt = \left[ \sum\limits_{j=1}^{n} P_{j}\dot{Q}_{j}-K \right]  dt \]




Resulta que esta $F$ que ha aparecido desde nuestro primer capítulo, resulta ser la función generatriz de la función canónica.



Tenemos:




\[ dF_{1} = \sum\limits_{j=1}^{n} \left[ P_{j} dq_{j}  - Hdt \right] - \sum\limits_{j=1}^{n} \left[ P_{j}dQ_{j}  - Kdt \right]  \]



Tenemos:


Significa entonces que:


\[ dF_{1}(q_{j},Q_{j},t) = \sum\limits_{j=1}^{n} \left[p_{j}dq_{j}-P_{j}dQ_{j}      \right] + \left[ K-H   \right]dt \]


Tenemos:

\[ \frac{ \partial F_{1} }{\partial q_{j}} = p_{j}  \]


\[  \frac{\partial F_{1}}{\partial Q_{j}} = - P_{j}; \quad \frac{\partial F_{1}}{\partial t} = K-H \quad \Rightarrow \quad \boxed{ K=H+ \frac{\partial F_{1}}{\partial t}  } \]


Tenemos:



\[ \dot{Q}_{i} = \frac{\partial K}{\partial P_{i}}  \]


\[ \dot{\partial}_{i} = - \frac{\partial K}{\partial Q_{i}}  \]

\textbf{Partícula libre}

\textbf{Oscilador armónico}


Para que $df_{1}(q_{j},Q_{j},t)$ sea una diferencial exacta se requiere que:



Tenemos:


\[ f(x,y) \]



\[ df = \frac{\partial f}{\partial x} dx + \frac{\partial f}{\partial y} dy \]



Si tenemos:



\[  \frac{\partial^{2}f}{\partial y \partial x} = \frac{ \partial^{2}f }{\partial x \partial y}  \]



\[ f(x,y,z) \]



\[ df = \frac{ \partial f }{\partial x} dx + \frac{\partial f}{\partial y} dy + \frac{\partial f}{\partial z}dz \]



Tenemos:



\[ dW(x,y,z) \overbrace{=}^{?} N(x,y,z)dx + \mu (x,y,z)dy + F(x,y,z) dz  \]



Tenemos:


\[ \frac{\partial F_{1}}{\partial q_{j}} = p_{j} \]


\[  \frac{\partial F_{1}}{\partial Q_{j}} = -P_{j} \]




Vamos a mostrar que el hamiltoniano es el generador de traslaciones temporales, se puede ver como una transformación infinitesimal, dentro de las transformaciones canónicas existen transformaciones, definiciones, es decir que analizando al hamiltoniano es posible que podamos resolver el problema.


Entonces podemos tener nosotros diferentes tipos de funciones generatrices, tenemos nosotros.



\[ F_{2} =F_{2}(q_{i},P_{i},t) = F_{1} + \sum\limits_{i=1}^{n}P_{i}Q_{i} \]



Tenemos:



\[  dF_{2} = p_{i}dq_{i} + Q_{i}dP_{i} + \left[  K-H  \right]dt \]



Tenemos:



\[ \frac{ \partial F_{2}  }{\partial q_{i}} = p_{i} \quad \frac{\partial F_{2}}{\partial t} = K-H \]

\[  \frac{\partial F_{2}}{\partial P_{i}} = Q_{i} \quad \boxed{ k=H+\frac{\partial F_{2}}{\partial t}  } \]



\begin{figure}[H]
\centering
    \includegraphics[width=0.7\textwidth]{1.png}
\caption*{.}
\end{figure}


Tenemos:


\[  (q_{i},p_{i}) \]



\[ P_{j} \equiv \frac{\partial L}{\partial \dot{q_{j}}} \]


Tenemos:


\[ (s_{i},\bar{p_{i}}) \]


Tenemos:


\[ \bar{p_{i}} = \frac{\partial L}{\partial \dot{s_{i}}}  \]



Resulta que la transformación de punto induce en el otro espacio una tranformación canónica en el espacio fase. 


Esas son de las transformaciones canónicas no genuinas, no triviales, esas vienen de una transformación de punto que se realizó desde el espacio de configuraciones. 


SI realizamos una transformación de Legendre, tenemos:


\[ H(q_{i},p_{i},t)  \to \quad T.L \]



EL sistema hamiltoniano que obtenemos es distinto del cuaj tenemos anteriormente.  


A nivel Lagrangiano no podemos mediante una tranformación de punto, si cambiamos el tiempo en partícula libre, podemos realizar la transformación canónica, siempre y cuando se cumplan los paréntesis de Poisson, eso es muy importante. 



Entonces ya que tenemos escrito lo anterior, ahora veremos lo siguiente respecto a las notas de profesor Ortigoza, este ejemolo nos dice a nosotros lo siguiente, supongamos que tenemos la siguiente transformación.



Tenemos:



\[ x(Q,P,t) = \frac{ P }{m \omega} \sen \left( \right)  \]



\begin{figure}[H]
\centering
    \includegraphics[width=0.7\textwidth]{2.png}
\caption*{.}
\end{figure}

Tenemos:


\[ K = \frac{P^{2}}{2m} \]

Tenemos:


\[ \dot{Q} = \frac{\partial K}{\partial P} = \frac{2P}{2m} = \frac{P}{m}  \]


Tenemos:


\[ \dot{P} = - \frac{\partial K}{\partial Q} = 0 \quad \to P=P_{0} = cte  \]



Tenemos:


\[ \frac{dQ}{dt} = \frac{ P_{0} }{m} \to \quad Q(t) = Q_{0} + \frac{P_{0}}{m}t   \]


Veremos otros ejemplo, tenemos:



\[ F_{2} = \sum\limits_{i=1}^{n} q_{i} P_{i}  \]


Tenemos:


\[ p_{k} = \frac{\partial F_{2}}{\partial q_{k}} = \frac{\partial }{\partial q_{k}} \left[ \sum\limits_{i=1}^{n} q_{i}P_{i}  \right] \]


\[  = \sum\limits_{i=1}^{n} \frac{ \partial q_{i}  }{\partial q_{k}} P_{i} = \sum\limits_{i=1}^{n} \delta_{ik}P_{i} = P_{k} \]


Tenemos:



\[  Q_{k} = \frac{\partial F_{2}}{\partial P_{k}} = \frac{ \partial  }{\partial P_{k}} \left[  \sum\limits_{i=1}^{n} q_{i}P_{i}  \right] = \sum\limits_{i=1}^{n}q_{i} \frac{\partial P_{i}}{\partial P_{k}} = \sum\limits q_{i}\delta_{ik} = q_{k} \]



Tenemos:



\[  K = H + \frac{\partial F_{2}}{\partial t}  = H \]


Tenemos:


\[ Q_{k}(q_{i},p_{i},t)  = q_{k} \]



\[ P_{k}(q_{i},p_{i},t) = p_{k}  \]



\begin{figure}[H]
\centering
    \includegraphics[width=0.7\textwidth]{3.png}
\caption*{.}
\end{figure}

\subsection{Transformaciones canónicas infinitesimales}


Partimos de una función de tipo $F_{2}$ que depende de las antiguas coordenadas, nuevos momentos y del tiempo.


Tenemos:



\[  (F_{2}(q_{i},P_{i},t))= \sum\limits_{i=1}^{n} q_{i}P_{i} + \epsilon G(q_{i},P_{i},t)  \]



\[ p_{k} = \frac{\partial F_{2}}{\partial q_{k}} = \frac{\partial}{\partial q_{k}} \left[ \sum\limits_{i=1}^{n} q_{i}P_{i} +\epsilon G(q_{i},P_{i},t)  \right]   \]


Tenemos:

\[ \boxed{ p_{k} = P_{k} + \epsilon \frac{\partial G}{\partial q_{k}} (q_{i},P_{i},t)  }    \]

Tenemos:



\[ Q_{k} = \frac{\partial F_{2}}{\partial P_{k}} = q_{k} + \epsilon \frac{\partial G}{\partial P_{k}} (q_{i},P_{i},t)  \]



Tenemos:



\[ K(Q_{i},P_{i},t) = H(q_{i},P_{i},t) + \frac{\partial F_{2}}{\partial t} \]


Tenemos: 



\[ K = H + \epsilon \frac{\partial G}{\partial t} \]


\begin{figure}[H]
\centering
    \includegraphics[width=0.9\textwidth]{4.png}
\caption*{.}
\end{figure}

Tenemos:



\[ \delta{q_{i}} = \epsilon \frac{ \partial G }{\partial p_{i}}  \]


\[ \delta p_{i} = -\epsilon \frac{\partial G}{\partial q_{i}}     \]



\[ \boxed{\frac{ \delta q_{i} }{\epsilon}    = \frac{\partial G}{\partial p_{i}}}  \]


\[ \boxed{\frac{\delta p_{i}}{\epsilon} = - \frac{\partial G}{\partial q_{i}}}  \]


Tienen la misma forma a estas:


\[ \frac{dq_{i}}{dt} = \frac{\partial H}{\partial p_{i}}   \]

\[ \frac{dp_{i}}{dt} = - \frac{\partial H}{\partial q_{i}}   \]

Tenemos:


\[  F_{2} = q_{i}P_{i}+\epsilon G \]


En realidad nos falta demostrar que esto:



\[ \delta{q_{i}} = \epsilon \frac{ \partial G }{\partial p_{i}}  \]


\[ \delta p_{i} = -\epsilon \frac{\partial G}{\partial q_{i}}     \]



Efectivamente cumple las propiedades de trnasformación canónica.


Entonces tenemos nosotros:



\[ Q_{i} = q_{i} + \epsilon \frac{\partial G}{\partial p_{i}} \]


\[ P_{i} = p_{i} - \epsilon \frac{\partial G}{\partial q_{i}}     \]


Tenemos:


\[ \left[  Q_{i},Q_{j}\right] = 0; \quad \left[ P_{i},P_{j}  \right] = 0; \quad \left[ Q_{i}, P_{j}   \right] =\delta{ij}    \]


Tenemos:

\[  \left[  Q_{i},Q_{j}  \right]  = \left[  q_{i} + \epsilon \frac{\partial G}{\partial p_{i}}, q_{j}+\epsilon \frac{\partial G}{\partial p_{j}}  \right]   \]


\begin{figure}[H]
\centering
    \includegraphics[width=0.9\textwidth]{5.png}
\caption*{.}
\end{figure}


\end{document}